\section*{Hoofdstuk \ref{ch:g9:acids-bases-ph} --- Zuren, basen en pH}

\begin{solution}{exo:g9:acids-bases-ph:1}
\omterm{def:g9:acids-bases-ph:ph}{pH} 3: zuur. \omterm{def:g9:acids-bases-ph:ph}{pH} 7: neutraal. \omterm{def:g9:acids-bases-ph:ph}{pH} 11: basisch.
\end{solution}

\begin{solution}{exo:g9:acids-bases-ph:2}
Zuur: \ce{H+}-ionen. Basisch: \ce{OH-}-ionen.
\end{solution}

\begin{solution}{exo:g9:acids-bases-ph:3}
Limoensap (2), koffie (5), zuiver water (7), zeewater (8.1),
huishoudammonia (12).
\end{solution}

\begin{solution}{exo:g9:acids-bases-ph:4}
Je legt een strookje op een schoon, droog schaaltje en brengt er met een
schoon glazen staafje een druppel van de oplossing op; de kleur
vergelijk je meteen met de kaart.
\end{solution}

\begin{solution}{exo:g9:acids-bases-ph:5}
Ongeveer 5 na een tienvoudige verdunning; ongeveer 6 na een honderdvoudige.
\end{solution}

\begin{solution}{exo:g9:acids-bases-ph:6}
Nee. Verdunnen brengt de \omterm{def:g9:acids-bases-ph:ph}{pH} dichter bij 7, maar nooit voorbij 7: water
zelf is neutraal, en door het toe te voegen kunnen de \ce{OH-}-ionen
nooit talrijker worden dan de \ce{H+}-ionen.
\end{solution}

\begin{solution}{exo:g9:acids-bases-ph:7}
Bij het \omterm{def:g2:mixing-and-dissolving:mix}{mengen} komt warmte vrij. Water dat op een geconcentreerd zuur
gegoten wordt, kan meteen koken en zuur wegslingeren; zuur dat langzaam
in water gegoten wordt, verspreidt zich meteen over een groot volume.
\end{solution}

\begin{solution}{exo:g9:acids-bases-ph:8}
Bijtend: tast huid, ogen en metalen aan. Draag een veiligheidsbril en
handschoenen; spoel elke spat weg met veel water.
\end{solution}

\begin{solution}{exo:g9:acids-bases-ph:9}
Drie verdunningen (1 naar 2, 2 naar 3, 3 naar 4): $10 \times 10 \times 10 =
1000$ keer.
\end{solution}

\begin{solution}{exo:g9:acids-bases-ph:10}
Magnesiamelk is basisch (\omterm{def:g9:acids-bases-ph:ph}{pH} ongeveer 10.5): het reageert met een deel
van het overtollige zuur in de maag.
\end{solution}

\begin{solution}{exo:g9:acids-bases-ph:11}
Nee: met een \omterm{def:g9:acids-bases-ph:ph}{pH} van 8.1 is hij nog steeds een beetje basisch. Hij gaat
richting een lagere \omterm{def:g9:acids-bases-ph:ph}{pH}, minder basisch, omdat er koolstofdioxide uit de
\omterm{def:g4:air-a-mixture-of-gases:air}{lucht} in \omterm{def:g2:mixing-and-dissolving:dissolve}{oplost}.
\end{solution}

\begin{solution}{exo:g9:acids-bases-ph:12}
\omterm{def:g9:acids-bases-ph:ph-paper}{pH-papier} lees je met het oog af tegen een kleurenkaart, tot op ongeveer
één eenheid; de \omterm{def:g9:acids-bases-ph:ph-paper}{pH-meter} geeft een getal op één of twee decimalen. De
meter is nauwkeuriger; de twee resultaten komen overeen binnen de
nauwkeurigheid van het papier.
\end{solution}

\begin{solution}{pb:g9:acids-bases-ph:1}
\textbf{1.} Zuur: het bevat vooral \ce{H+(aq)}- en \ce{Cl-(aq)}-ionen.

\textbf{2.} GHS05, bijtend: een zuur tast huid, ogen en metalen aan.

\textbf{3.} Hij geeft een nauwkeurig getal in plaats van een kleur die je
met het oog moet vergelijken.

\textbf{4.} Tien keer meer: één pH-eenheid komt overeen met een
tienvoudige verdunning.

\textbf{5.} \qty{10}{L}, met \omterm{def:g9:acids-bases-ph:ph}{pH} 3.

\textbf{6.} Drie: \omterm{def:g9:acids-bases-ph:ph}{pH} 2 naar 3, 3 naar 4, 4 naar 5.

\textbf{7.} $1 \times 10 \times 10 \times 10 = \qty{1000}{L}$.

\textbf{8.} Nee: verdunnen brengt de \omterm{def:g9:acids-bases-ph:ph}{pH} alleen dichter bij 7.

\textbf{9.} \ce{H+ + OH- -> H2O}.

\textbf{10.} Het zuur wordt vernietigd, omgezet in water, in plaats van
over een enorm volume verspreid; er is maar een beetje base nodig en
geen duizend liter water.

\textbf{11.} Baksoda opgelost in water (of zeepsop).

\textbf{12.} $1000 - 1 = \qty{999}{L}$ water.
\end{solution}
